% Generated draft: requires mathematical and source review.
\documentclass[11pt]{article}
\usepackage[margin=25mm]{geometry}
\usepackage{fontspec}
\setmainfont{Latin Modern Roman}
\usepackage{amsmath,amssymb,mathtools,unicode-math}
\setmathfont{Latin Modern Math}
\usepackage{enumitem,needspace,xurl,hyperref}
\setlength{\parindent}{0pt}
\setlength{\parskip}{5pt}
\setlength{\emergencystretch}{3em}
\newcommand{\N}{\mathbb{N}}
\newcommand{\Z}{\mathbb{Z}}
\newcommand{\Q}{\mathbb{Q}}
\newcommand{\R}{\mathbb{R}}
\newcommand{\C}{\mathbb{C}}
\newcommand{\F}{\mathbb{F}}
\newcommand{\A}{\mathbb{A}}
\newcommand{\PP}{\mathbb{P}}
\newcommand{\E}{\mathbb{E}}
\newcommand{\eps}{\varepsilon}
\newcommand{\dd}{\,\mathrm{d}}
\begin{document}
\section*{UnsolvedMath ID: 800579}
\textbf{Source problem code:} catalogue entry
\subsection{Definitions and mathematical statement}
Regard a set of reals \(M\subseteq\R\) as a space with the usual subspace topology, and put \(\N=\{1,2,\ldots\}\). The supplied excerpt does not define its cover conventions; the following gives the standard topological interpretation, subject to confirmation against the source.

An open cover of \(M\) is a family \(\mathcal U\) of open subsets of \(M\) such that \(\bigcup\mathcal U=M\) and \(M\notin\mathcal U\). Write \(\mathcal O(M)\) for the class of these covers. An \(\omega\)-cover is an open cover \(\mathcal U\) such that every finite subset of \(M\) is contained in some member of \(\mathcal U\); write \(\Omega(M)\) for this class. A \(\gamma\)-cover is an infinite open cover \(\mathcal U\) such that, for every \(x\in M\), the set \(\{U\in\mathcal U:x\notin U\}\) is finite; write \(\Gamma(M)\) for this class. These definitions impose no countability restriction on individual covers.

The property \(\mathsf S_{\mathrm{fin}}(\Omega,\Omega)\) means that for every sequence \((\mathcal U_n)_{n\in\N}\) in \(\Omega(M)\), there are finite subfamilies \(\mathcal V_n\subseteq\mathcal U_n\) such that
\[
\bigcup_{n\in\N}\mathcal V_n\in\Omega(M).
\]
The property \(\mathsf U_{\mathrm{fin}}(\mathcal O,\Gamma)\) means that for every sequence \((\mathcal U_n)_{n\in\N}\) in \(\mathcal O(M)\), each having no finite subcover of \(M\), there are finite subfamilies \(\mathcal V_n\subseteq\mathcal U_n\) such that
\[
\left\{\bigcup\mathcal V_n:n\in\N\right\}\in\Gamma(M).
\]
Thus failure of this latter property means that some such sequence admits no sequence of finite selections producing a \(\gamma\)-cover in this way.

The source asks: Is there a direct, non-dichotomic construction of a set \(M\subseteq\R\) satisfying
\[
\mathsf S_{\mathrm{fin}}(\Omega,\Omega)
\quad\text{and}\quad
\neg\mathsf U_{\mathrm{fin}}(\mathcal O,\Gamma)?
\]
The requested construction method is part of the question. The supplied record does not define ``direct'' or ``non-dichotomic'', identify the dichotomy to be avoided, or specify any additional set-theoretic hypotheses. Consequently, the displayed properties specify the requested set under the conventions above, but do not fully formalize the methodological requirement.
\subsection{Short English statement}
Find a direct, non-dichotomic construction of a set of reals with \(\mathsf S_{\mathrm{fin}}(\Omega,\Omega)\) but without \(\mathsf U_{\mathrm{fin}}(\mathcal O,\Gamma)\). The meaning of the construction requirement and the source's exact cover conventions require confirmation. The registry supplies an open label; no independent status verification is asserted.
\subsection{Sources}
\begin{itemize}
\item Supplied local UnsolvedMath registry record: permanent ID \(800579\), source problem code \texttt{catalogue entry}, title ``Selection Principles in Mathematics --- Problem 5 ([SFH])''. The record lists Tsaban as proposer and \(2006\) as the proposed year; its supplied status is \texttt{open}.
\item Tsaban, \emph{Selection Principles and Special Sets of Reals: Open Problems} (2006), Problem 5. Supplied source URL: \url{https://arxiv.org/abs/math/0609601}. The record reports extraction by \texttt{arxiv\_latex} and an access date of \texttt{2026-10-16}; these are supplied provenance details, not a claim of independent source access or verification.
\item The supplied background gives attribution/context as \texttt{[SFH]}; its expansion and bibliographic referent are not supplied.
\end{itemize}
\end{document}
